\section{模型检验与敏感性}

本节不另开求解器，只把已经冻结的开发集、配对种子、留出集与寻路版本串成检验。口径一律标为非正式或离线；正式测试仍然缺失。检验的目的不是补一张正式成绩表，而是回答四件事：问题1的覆盖判定会不会写死；问题2的 $J(q)$ 会不会签发虚假证书；问题3把预算写成停机会不会破坏全清；问题4换格点会不会在漏清未修时假装变快。

\subsection{问题1的几何不变量}

等边三角形演示把直径圆覆盖判定钉在否上。
直径约为 $1.0000000000134766$，Jung 半径约为 $0.5773502691948129$，第三顶点到直径中点的距离为 $\sqrt{3}/2>1/2$。
两点楔形演示则给出直径 $9.877085164818205$、包围半径 $4.938542582409102$ 的四边形，此时直径圆恰好覆盖，说明算法会随顶点集合切换结论，而不是写死某一个布尔值。
空集、单点、无界楔在实现里各有返回状态；开发脚本未把无界情形写进演示数据，正文不以虚构数字补全。
数值容差 $\tau=10^{-8}$ 只用于顶点过滤，不改变等边反例的几何判定。
后问使用的 $1.005^\circ$ 数值包络不进入问题1的演示，以免把量化约定冒充题面误差。

\subsection{问题2开发集}

分层种子 $n=20$ 上，一级不变量全部通过。
真值落入第一层外包失败 $0$，落入第二层后验失败 $0$，外包包含真源失败 $0$，核心层虚假保证 $0$，虚假清除证书 $0$，同半径约束违反 $0$。
失败种子列表为空。
无信号支在二十例上为 $0$，因此开发集不能用来讨论第二点收不到射频的情形，鲁棒性里的“无信号支不得挖球”只能作为实现断言，不能写成该批经验频率。
$J^\ast$ 的中位数 $42.75989\,\mathrm{m}$，$90\%$ 分位 $45.40667\,\mathrm{m}$，$95\%$ 分位 $45.48670\,\mathrm{m}$，最大 $46.20539\,\mathrm{m}$，全部明显高于 $20\,\mathrm{m}$。
真值后验半径的中位数 $16.041\,\mathrm{m}$，$90\%$ 分位 $33.988\,\mathrm{m}$，最大 $44.240\,\mathrm{m}$；后验半径不超过 $20\,\mathrm{m}$ 的比例 $13/20$，保证第二点比例 $7/20$。
二十次第二反馈均为 $\mathtt{direction}$，没有 $\mathtt{near}$ 或 $\mathtt{no\_signal}$。
该生成器把第一点放在可测位置，因而不能外推到任意第一点都收得到信号。
行走时间中位数约 $202.8318\,\mathrm{s}$，$90\%$ 分位约 $214.37\,\mathrm{s}$，最大约 $270.40\,\mathrm{s}$，与演示的 $191.74\,\mathrm{s}$ 同量级。
候选点数中位数 $62$，与演示的 $191$ 不同，因为演示含局部加密层。
$J(q^\ast)$ 相对在 $c_1$ 处直接取值的均值差为 $-271.014$，说明排序确实把候选从包围圆心推开，而不是退化为走到 $c_1$。
分层按方位八层：第一层五例中位数 $44.579\,\mathrm{m}$，第二层三例 $42.798\,\mathrm{m}$，第三层两例 $42.693\,\mathrm{m}$，第四层两例 $43.912\,\mathrm{m}$，第五层两例 $26.222\,\mathrm{m}$，第六层两例 $43.620\,\mathrm{m}$，第七层两例 $42.698\,\mathrm{m}$，第八层两例 $41.670\,\mathrm{m}$。
第五层较低并不改变全体高于 $20\,\mathrm{m}$ 的结论，也不能把某一层写成已经可清。
相对加密参考集的 $J$ 间隙中位数约 $0.055\,\mathrm{m}$，最大约 $0.348\,\mathrm{m}$，只说明分层候选没有相对加密网格大幅变差。
示向度微扰抖动在该批上样本数为 $0$，不把空统计写成鲁棒性结论。
候选点数中位数 $62$，最大 $63$；CPU 中位数约 $0.235\,\mathrm{s}$，最大约 $0.269\,\mathrm{s}$。
真值后验半径 $95\%$ 分位 $34.528\,\mathrm{m}$ 介于 $90\%$ 分位与最大值之间，与 $13/20$ 的经验频率同一口径。
声明限制为自建分层场景，不是竞赛概率，也不是连续全局最优。

\subsection{问题3策略消融}

同晚十局把三条策略放在同一演练入口下对照，但局间场景独立，不是配对种子。
GREEDY\_FAST 相对 GREEDY 的平均 $T/K$ 由 $390\,\mathrm{s}$ 降到 $341\,\mathrm{s}$，SAFE 均值由 $1.20$ 降到 $0.20$，全清率保持 $10/10$。
ABORT 把源预算写成放弃条件后，全清率掉到 $1/10$，平均 $T/K=382\,\mathrm{s}$ 并不能补偿漏清。
该对照的含义是完备性对字典序目标不可谈判：任何把启发式预算升级为停机规则的改动，都应在全清率上被否决，而不是只看平均时间。
十局 $T/K$ 的区间宽度约 $200\,\mathrm{s}$（主线 $244$--$450$），反映源的几何散布与发现顺序，而不是随机数种子可重复的置信区间。
另十局轨迹均值 $340.98\,\mathrm{s}$，区间 $[243.74,449.50]$，与取整后的 $341\,\mathrm{s}$ 一致。
预算模式 $\mathrm{SourceBudgetS}=300$ 不作为主线。
GREEDY 快照保留为回退基线，正式入口默认 GREEDY\_FAST。
口径全部为非正式演练，禁止填入表1。
ABORT 的 $1/10$ 已经足够否决该候选：字典序下全清率优先，不必再讨论它的平均时间是否碰巧接近 GREEDY。
GREEDY\_FAST 的区间下沿 $244\,\mathrm{s}$ 低于 GREEDY 的 $334\,\mathrm{s}$，上沿 $450\,\mathrm{s}$ 与 GREEDY 的 $455\,\mathrm{s}$ 接近，说明加速主要发生在走廊较少的局，而不是把所有局都削去同一常数。

\subsection{问题4求解器修复与留出}

SQUARE81 在种子 $0$--$29$ 上修复前后对照如下。
全清且证书由 $13/30$ 升至 $30/30$，平均 $T/K$ 由 $2298.44\,\mathrm{s}$ 降至 $1313.16\,\mathrm{s}$，相对下降 $42.87\%$。
平均 $T$ 由 $28875.13\,\mathrm{s}$ 降至 $17844.98\,\mathrm{s}$，下降 $38.20\%$；平均路程由 $120722.01\,\mathrm{m}$ 降至 $68250.75\,\mathrm{m}$，下降 $43.46\%$。
旧版部分场景漏清，效率比较必须与全清结果一起报告，否则会把中途退出的短时间误读成更快。
旧版完整演练 $13/16$、$7/11$、$9/13$ 的 $T/K$ 为 $3025.96$、$5593.18$、$4427.19\,\mathrm{s}$，路程 $169.46$、$158.00$、$164.26\,\mathrm{km}$，与 $81$ 点网格约 $48\,\mathrm{km}$ 的基本行程对照，说明高耗时首先来自往返调度而不是格点本身。

留出种子 $1000$--$1029$、误差固定取 $\pm 1^\circ$ 两端再量化，全清且证书仍为 $30/30$，平均 $T/K=1376.27\,\mathrm{s}$。
留出集略慢于修复后的训练种子，但没有漏清。
该模式专门打击测试时方位无噪声的盲区，不能解释成对官方模拟器误差实现的统计覆盖。
离线单场墙钟由约 $11\,\mathrm{s}$ 降至 $0.1$--$0.4\,\mathrm{s}$，使三十种子配对成为可重复实验。

\subsection{覆盖层与寻路}

同一修复后求解器、同一批种子上，HEX37 相对 SQUARE81 的 $T/K$ 由 $1313.160522935798\,\mathrm{s}$ 降至 $836.9129462514154\,\mathrm{s}$，绝对下降 $476.2475766843827\,\mathrm{s}$，相对下降 $36.267\%$。
平均 $T$ 下降 $36.019\%$，路程下降 $32.314\%$，测向次数下降 $50.407\%$，覆盖访问由 $81$ 降至 $37$（$54.321\%$）。
光学兜底次数几乎不变，约 $1.067$ 对 $1.167$，相对变化为略增；清除动作由 $60.567$ 降至 $58.5$，相对下降约 $3.412\%$，说明减少的是发现层，而不是把清除几何变松。
配对下降量在三十个种子上均为正，没有个别种子上 HEX37 更慢且漏清的反例。
$N=16$ 种子七个，当前仍走完覆盖点证明其余四频道缺席。

寻路版本在 HEX37 上离线均值 $836.91\to 774.06\to 735.31\to 711.20\to 683.09$，全程 $30/30$ 证书。
V1 把插入阈值收到 $400\,\mathrm{m}$ 时路程由约 $44.2\,\mathrm{km}$ 降至 $36.1\,\mathrm{km}$，但证书模式把测向次数从 $337$ 抬到 $545$。
V2 把路径上待清除批处理后测向回到 $375$，路程回到 $38.3\,\mathrm{km}$，$T/K$ 由 $774$ 降至 $735$。
V3 冻结清除插入上限 $600\,\mathrm{m}$、测向插入上限 $250\,\mathrm{m}$、激进试清半径 $80\,\mathrm{m}$，再降 $3.3\%$ 到 $711.20$。
PREFIX\_A 相对 V3 再降约 $4\%$，第一次检出序号 $6.85\to 4.20$。
阈值继续加密的收益已经小于百分之五，停止调参，以免把离线生成器拟合进插入常数。
非正式演练 PREFIX\_A 两场为 $11/11$、$822.38\,\mathrm{s}$ 与 $10/10$、$930.28\,\mathrm{s}$；更早 HEX37 两场 $12/12$、$899.04$ 与 $16/16$、$620.08$；SQUARE81 两场 $12/12$、$1654.84$ 与 $16/16$、$1032.43$。
场景非配对，不能单独归因百分比。

\subsection{检验边界}

以上数字有三条共同边界。
离线几何模型复现题设计时与 $\pm 1^\circ$ 量化误差，但不经过接口、不经过模拟器的墙钟与登录校验。
演练测试经过接口，但场景随机非配对，局数少，不能给出竞赛概率。
正式测试未执行，任何把上述均值写入表1的做法都会把口径弄错。
回归测试 $168$ 项通过，只说明当前实现满足仓库内断言，不说明一切未写入断言的场景都已穷尽。
HEX37 真子集是否仍覆盖一切朝向尚未另证；问题3的七点对 $r_i\in[1000,1200)$ 不是单点可测的严格证书。
截止时间之后不能补做正式测试。

把四问检验按不变量、消融、配对、留出排列，是为了防止用同一批数字回答不同问题。
问题1只问覆盖判定会不会写死，不需要三十种子。
问题2只问 $J$ 会不会签发虚假证书，开发集 $n=20$ 已经足够暴露 $J^\ast$ 全体高于 $20\,\mathrm{m}$。
问题3的对照回答的是停机规则，不是期望 $T/K$。
问题4必须先看全清率，再看时间：修复前 $13/30$ 的均值 $2298.44\,\mathrm{s}$ 含漏清，不能单独拿来和修复后的 $1313.16\,\mathrm{s}$ 比快慢而不提 $30/30$。
HEX37 相对 SQUARE81 的 $36.267\%$ 只在求解器已修复、种子配对的前提下成立。
PREFIX\_A 的 $683.09\,\mathrm{s}$ 是离线寻路均值，与演练 $822.38\,\mathrm{s}$、$930.28\,\mathrm{s}$ 不是同一口径，禁止平均成一个现场成绩。
寻路链 $836.91\to 774.06\to 735.31\to 711.20\to 683.09$ 全程 $30/30$，阈值继续加密的收益已经小于百分之五，停止调参。
V1 在插入阈值 $400\,\mathrm{m}$ 时路程由约 $44.2\,\mathrm{km}$ 降至 $36.1\,\mathrm{km}$，但测向从 $337$ 抬到 $545$；V2 批处理后测向回到 $375$，路程回到 $38.3\,\mathrm{km}$。
这些中间版本说明“更短路程”与“更少测向”可以互相顶牛，必须同时报告全清率，不能只看一项。
留出端点误差均值 $1376.27\,\mathrm{s}$ 略慢于训练种子修复后的 $1313.16\,\mathrm{s}$，但没有漏清，不能把留出变慢写成过拟合失败，也不能写成正式成绩上界。
